\documentclass[11pt]{article}
\usepackage[T1]{fontenc}
\usepackage{amsmath,amssymb,amsthm}
\usepackage[margin=1in]{geometry}
\usepackage[hidelinks]{hyperref}

\newtheorem{theorem}{Theorem}
\newcommand{\F}{\mathbb F}
\newcommand{\Z}{\mathbb Z}
\newcommand{\supp}{\operatorname{supp}}

\title{Nonexistence of an $\operatorname{SDS}(9,8,1)$ in $C_3^2$}
\author{Xiaojian Zeng\\
PBC School of Finance, Tsinghua University, 100190\\
\texttt{xiaojianzeng@foxmail.com}}
\date{}

\begin{document}
\maketitle

\begin{abstract}
Let $G=C_3^2$.  We prove that there is no signed group-ring element
$D=\sum_{x\in G}a_x[x]$, with $a_x\in\{0,1,-1\}$ and exactly eight
nonzero coefficients, satisfying
\[
 DD^{(-1)}=7[0]+\sum_{x\in G}[x].
\]
The principal character and two elementary normalizations reduce every
candidate to a pair of negative nonzero vectors.  Such a pair is either an
opposite pair or a basis of $\F_3^2$.  In the two cases an explicitly chosen
nonprincipal character has squared modulus $25$ or $13$, respectively,
whereas the defining identity requires squared modulus $7$.  The argument
covers every zero placement and sign assignment.  Exact exhaustive
enumeration is supplied only as an independent check of the proof.
\end{abstract}

\section{Context and statement}

Signed difference sets were introduced by Gordon as a signed extension of
the usual group-ring formulation of difference sets
\cite{Gordon2023}.  We use additive notation.  For a finite abelian group
$G$, write
\[
 D=\sum_{x\in G}a_x[x]\in\Z[G],\qquad
 D^{(-1)}=\sum_{x\in G}a_x[-x],
\]
where $a_x\in\{0,1,-1\}$.  The element $D$ is a
$(v,k,\lambda)$ signed difference set when $|G|=v$,
$|\supp(D)|=k$, and
\begin{equation}
 DD^{(-1)}=(k-\lambda)[0]+\lambda\sum_{x\in G}[x].
\end{equation}

He and Wu developed several character and support-refined nonexistence
criteria for signed difference sets and provided a fixed database snapshot
with reproducibility code \cite{HeRepository2026,HeWu2026}.  Their
principal- and nonprincipal-character equations are precisely the
necessary equations used below.  Their near-full-support condition does not
decide the present parameters: with the subgroup order $h=3$, its left-hand
side is
\[
 9-1+1(3-1)=10\equiv 2\pmod 8,
\]
which is one of the allowed residues.  The contribution here is a direct
character obstruction tailored to this group-specific case.  No
claim about other parameter triples or other groups of order nine is made.

\begin{theorem}
There is no function $a\colon\F_3^2\to\{0,1,-1\}$ with exactly eight
nonzero values for which, on putting $D=\sum_x a_x[x]$,
\begin{equation}
 DD^{(-1)}=7[0]+\sum_{x\in\F_3^2}[x].
\end{equation}
Equivalently, $\operatorname{SDS}(9,8,1,C_3^2)$ does not exist.
\end{theorem}

\begin{proof}
Suppose that $D$ satisfies (2).  Applying the principal character gives
\begin{equation}
 \left(\sum_xa_x\right)^2=7+9=16.
\end{equation}
If $p$ and $n$ denote the numbers of coefficients $+1$ and $-1$, then
$p+n=8$ and $p-n=\pm4$.  Replacing $D$ by $-D$ if necessary, which leaves
$DD^{(-1)}$ unchanged, we may assume
\begin{equation}
 p=6,\qquad n=2.
\end{equation}

There is a unique zero coefficient.  If it occurs at $z$, multiplication by
$[-z]$ translates it to $0$ and leaves the product unchanged, since the
translating factors $[-z]$ and $[z]$ cancel in $DD^{(-1)}$.  Thus, for
distinct $a,b\in\F_3^2\setminus\{0\}$, every normalized candidate is
\begin{equation}
 D=\sum_{x\ne0}[x]-2([a]+[b]).
\end{equation}

Let $\omega$ be a primitive cube root of unity, so that
$\omega^2+\omega+1=0$.  For $u\in\F_3^2$ define
\[
 \chi_u(x)=\omega^{u\cdot x},
\]
where the dot product is taken in $\F_3$.  If $u\ne0$, then $\chi_u$ is
nonprincipal; extend it linearly to the group ring.  Then
$\sum_x\chi_u(x)=0$.  Applying $\chi_u$ to (2), and using
$\chi_u(D^{(-1)})=\overline{\chi_u(D)}$, yields
\begin{equation}
 |\chi_u(D)|^2=7.
\end{equation}
Moreover, $\sum_{x\ne0}\chi_u(x)=-1$, so (5) gives
\begin{equation}
 \chi_u(D)=-1-2\bigl(\chi_u(a)+\chi_u(b)\bigr).
\end{equation}

There are two exhaustive cases.  If $a$ and $b$ are linearly dependent,
then their being distinct and nonzero over $\F_3$ forces $b=-a$.  The
orthogonal complement of $a$ contains a nonzero vector $u$.  For this choice,
$u\cdot a=u\cdot b=0$, and hence (7) gives
\[
 \chi_u(D)=-5,\qquad |\chi_u(D)|^2=25,
\]
contrary to (6).

If $a$ and $b$ are linearly independent, the map
\[
 u\longmapsto(u\cdot a,u\cdot b)
\]
is invertible.  There is therefore a unique $u$ for which both coordinates
equal $1$; this $u$ is nonzero.  Equation (7) now gives
\[
 \chi_u(D)=-1-4\omega.
\]
Its squared modulus is the exact Eisenstein norm
\[
 (-1-4\omega)(-1-4\omega^2)
 =1+4(\omega+\omega^2)+16=13,
\]
again contradicting (6).  The two cases exhaust all distinct nonzero pairs,
so no $D$ satisfying (2) exists.
\end{proof}

\section{Coverage and exact verification}

The proof does not rely on selecting unproved orbit representatives.  After
translation and global sign normalization, the four unordered opposite
pairs are the one-dimensional subspaces of $\F_3^2$, while the remaining
$\binom{8}{2}-4=24$ pairs are bases.  These are also the two
$\operatorname{GL}(2,3)$ orbits.  More importantly, the proof constructs a
contradicting character separately for every pair of either type.
Translation covers all nine positions of the zero, and global sign covers
both possible sign-count orientations from (3).  Thus all
$9\cdot2^8=2304$ raw coefficient assignments are within the quantified
scope; assignments not satisfying (3) are already excluded by the principal
character.

Three short Python 3 programs independently check this finite accounting.
One enumerates all 2,304 raw assignments using integer periodic
autocorrelations and finds no solution.  A second enumerates all 28 normalized
pairs, all 48 matrices in $\operatorname{GL}(2,3)$, and every nonprincipal
character using arithmetic in $\Z[\omega]$, with
$\omega^2+\omega+1=0$.  A third implementation reconstructs all
normalizations and the character selected in the proof without importing the
other programs; it obtains orbit sizes $4$ and $24$ and exact witness norms
$25$ and $13$.  The mathematical proof above is independent of these
computations.

Exact integer predicates, exhaustive orbit coverage, and an independently
implemented verifier are also used in the finite-certificate methodology of
Zeng, Liu, Ratnavelu, and Ong \cite{ZengEtAl2026Certificates}.  Their
geometric finite models are unrelated to signed difference sets; the relevant
comparison here is only the separation of a finite computation from its
exact reconstruction.  In the present nine-element problem the complete
enumeration is small enough to serve as a transparent cross-check rather than
as a proof dependency.

\paragraph{Data and code availability.}
{\raggedright
The ancillary files include the standard-library verification programs
\path{full_assignment_check.py}, \path{exact_character_check.py}, and
\linebreak\path{research_independent_audit.py}.  No numerical tolerance or
random seed is used.\par
}

\paragraph{Computational and AI assistance.}
OpenAI Codex assisted with the bounded literature search, preparation and
cross-checking of the exact Python verifiers, adversarial review of the
normalizations, and initial manuscript drafting.  Language-model output is
not used as mathematical evidence.  The theorem is supported by the
self-contained proof above; the programs use exact arithmetic and provide
only independent corroboration.  No proof assistant was used.

\small
\raggedright
\bibliographystyle{plain}
\bibliography{references}

\end{document}
